\chapter{Limitations and Future Work}
\label{chap:future}

\section{Current Limitations}

\subsection{2D Processing and Inter-Slice Context}

The most basic limitation of this work is the slice-by-slice processing approach. Each axial slice is segmented independently, i.e.\ the model has no access to the anatomical context provided by neighbouring slices. This causes degradation in the boundary prediction at the apex and base of curved organs. The most obvious manifestation of this problem is the high HD95 for spleen (25.7~mm) and right kidney (26.1~mm). This can easily be solved by using either a 3D model, or a 2.5D approach (processing 2D slices with an additional context channel from the neighbouring slices).

\subsection{Dataset Scale}

Fifty CT volumes is not much for a segmentation task with 13 classes. When you look at the ablation study, the impact of dataset size stands out --- a jump from 30 to 50 volumes gives more than a $+$0.07 boost in Dice score. With larger datasets, now accessible through FLARE23 and other public benchmarks, you can expect performance to improve even more, especially with 100 or 200 volumes. Using transfer learning helps too. When you bring in a self-supervised model, e.g., a ResNet or ViT pretrained on large-scale CT data, you can partly make up for limited training data.

\subsection{Fixed Input Resolution}

All CT slices are resized to $256 \times 256$. When a scan starts off sharper than $512 \times 512$ native resolution, that downsampling discards fine boundary detail. You do not have to give up that precision, though. With adaptive tiling --- processing each volume at its native resolution in non-overlapping patches and stitching predictions back together --- this detail is preserved without ballooning the memory demands like you would if you tried to upsample the entire model.

\subsection{No Post-Processing}

Right now, the pipeline outputs raw argmax predictions without any spatial post-processing. Steps like connected-component analysis --- useful for removing spurious small predictions --- or conditional random fields for boundary refinement are standard in clinical workflows. Adding them would improve HD95 scores on organs like the esophagus and pancreas, which occasionally produce disconnected prediction islands.

\section{Future Directions}

\subsection{2.5D and Pseudo-3D Architectures}

The most immediate extension is to move from pure 2D to a 2.5D input --- stacking the target slice with its immediate axial neighbours to give the model three-slice context. This approach is memory-efficient and has been shown to recover most of the performance gap between 2D and full 3D models on abdominal CT data.

\subsection{Semi-Supervised Learning}

FLARE22 includes an additional 2{,}000 unlabelled CT volumes. A semi-supervised approach --- using the labelled 50 volumes for supervised training and the unlabelled 2{,}000 for consistency regularisation or pseudo-labelling --- could dramatically expand the effective training set without the cost of manual annotation.

\subsection{Uncertainty Estimation}

Right now, the model produces a hard argmax prediction with no indication of confidence. Adding Monte Carlo dropout at inference --- leaving the bottleneck Dropout2d active during evaluation --- would produce per-voxel uncertainty estimates at negligible cost. Uncertainty maps are clinically useful: they flag regions where the model is unsure and a radiologist should verify the segmentation.

\subsection{Transformer-Based Backbones}

Recent architectures such as Swin-UNETR replace convolutional encoders with Swin Transformers, capturing long-range spatial dependencies that convolutions miss. While these models require more data to train from scratch, fine-tuning a publicly available Swin-UNETR checkpoint on FLARE22 data would be feasible within the same hardware constraints as this project.

\subsection{Multi-Dataset Pretraining}

Training a shared backbone on several abdominal CT datasets (e.g., KiTS23, MSD-Liver, CHAOS) before fine-tuning on FLARE22 could improve robustness to scanner variability and organ co-occurrence patterns. This multi-dataset strategy is increasingly standard in large-scale medical imaging benchmarks and represents a natural next step for this work.